% 사물인터넷과 빅데이터 - 강의계획서
% AI디지털리터러시 문서 스타일과 동일하게 구성
% XeLaTeX 컴파일 필요
\documentclass[11pt, a4paper]{article}

\usepackage{fontspec}
\setmainfont{Noto Serif CJK KR}
\setsansfont{Noto Sans CJK KR}

\usepackage{geometry}
\geometry{a4paper, left=20mm, right=20mm, top=20mm, bottom=20mm}

\usepackage{xcolor}
\usepackage{tikz}
\usetikzlibrary{shapes.geometric, arrows.meta, positioning}
\usepackage{booktabs}
\usepackage{array}
\usepackage{longtable}
\usepackage{multirow}
\usepackage{enumitem}
\usepackage{fancyhdr}
\usepackage{titlesec}
\usepackage{hyperref}
\usepackage{colortbl}
\usepackage{tabularx}
\usepackage{makecell}
\usepackage{float}
\usepackage{tcolorbox}
\tcbuselibrary{skins, breakable}

% ========== 색상 정의 (밝은 톤) ==========
\definecolor{mainblue}{RGB}{41, 128, 185}
\definecolor{headerblue}{RGB}{214, 234, 248}
\definecolor{lightblue}{RGB}{235, 245, 251}
\definecolor{accentorange}{RGB}{230, 126, 34}
\definecolor{lightorange}{RGB}{253, 235, 208}
\definecolor{accentgreen}{RGB}{39, 174, 96}
\definecolor{lightgreen}{RGB}{212, 239, 223}
\definecolor{warningred}{RGB}{231, 76, 60}
\definecolor{lightred}{RGB}{250, 219, 216}
\definecolor{textgray}{RGB}{44, 62, 80}
\definecolor{lightgray}{RGB}{248, 249, 250}

\hypersetup{colorlinks=true, linkcolor=mainblue, urlcolor=mainblue}

% ========== 장/절 스타일 (페이지 브레이크 없음) ==========
\titleformat{\section}
    {\Large\bfseries\color{mainblue}}
    {\thesection.}
    {0.5em}
    {}
\titlespacing*{\section}{0pt}{15pt}{10pt}

\titleformat{\subsection}
    {\large\bfseries\color{mainblue}}
    {\thesubsection}
    {0.5em}
    {}
\titlespacing*{\subsection}{0pt}{12pt}{6pt}

\titleformat{\subsubsection}
    {\normalsize\bfseries\color{textgray}}
    {\thesubsubsection}
    {0.5em}
    {}
\titlespacing*{\subsubsection}{0pt}{8pt}{4pt}

% ========== 헤더/푸터 ==========
\pagestyle{fancy}
\fancyhf{}
\fancyfoot[C]{\thepage}
\renewcommand{\headrulewidth}{0pt}

% ========== 표/그림 번호 한글화 ==========
\renewcommand{\tablename}{표}
\renewcommand{\figurename}{그림}

% ========== 표 스타일 ==========
\renewcommand{\arraystretch}{1.4}
\newcolumntype{C}[1]{>{\centering\arraybackslash}p{#1}}
\newcolumntype{L}[1]{>{\raggedright\arraybackslash}p{#1}}

% ========== 인용문 스타일 ==========
\newcommand{\note}[1]{%
    \vspace{3pt}
    \noindent\textcolor{textgray}{\textbf{>} #1}
    \vspace{3pt}
}

% ========================================
\begin{document}

% ========== 제목 ==========
\begin{center}
{\fontsize{24pt}{28pt}\selectfont\bfseries\color{mainblue} 사물인터넷과 빅데이터}

\vspace{8pt}

{\Large\bfseries 강 의 계 획 서}

\vspace{8pt}

{\normalsize\bfseries 문서 버전: v1.5.4.0}
\end{center}

\vspace{5pt}
\noindent\rule{\textwidth}{0.5pt}

% ========================================
\section{교과목 기본 정보}
% ========================================

\begin{table}[H]
\centering
\begin{tabular}{|C{2.2cm}|L{4.5cm}|C{2.2cm}|L{4.5cm}|}
\hline
\rowcolor{headerblue}
\textbf{항목} & \textbf{내용} & \textbf{항목} & \textbf{내용} \\
\hline
\textbf{교과목명} & 사물인터넷과 빅데이터 & \textbf{학점/시간} & 2학점 / 주 2시간 \\
\hline
\textbf{수업 기간} & 15주 (총 30시간) & \textbf{수업 대상} & 컴퓨터교육과 4학년 \\
\hline
\textbf{수업 방식} & 프로젝트 중심 실습 & \textbf{수강 인원} & 21명 \\
\hline
\textbf{강의실} & 정306호 (컴퓨팅교육실습실) & \textbf{담당 교수} & 양창모 \\
\hline
\textbf{연락처} & cmyang@cje.ac.kr & \textbf{교재} & 담당교수 제공 자료 \\
\hline
\end{tabular}
\end{table}

\vspace{3pt}
\noindent\rule{\textwidth}{0.3pt}

% ========================================
\section{교과목 개요}
% ========================================

본 교과목은 마이크로비트와 Grove 확장 센서를 활용하여 데이터를 수집하고, 이를 분석·해석하는 과정을 통해 \textbf{컴퓨팅 사고력}, \textbf{인공지능 사고력}, \textbf{데이터 리터러시}를 함양하는 것을 목표로 한다. 학생들은 \textbf{학교 현장에 적용할 수 있는 프로젝트}를 직접 선정하고, 기획·제작·수집·분석의 전 과정을 수행한다. 이 과정에서 \textbf{바이브 코딩}(AI 협업 코딩)을 통해 코드 문법 없이 데이터 시각화, 분석, AI 모델 구현을 경험한다.

\vspace{3pt}
\noindent\rule{\textwidth}{0.3pt}

% ========================================
\section{핵심 메시지}
% ========================================

\begin{enumerate}[leftmargin=*, itemsep=3pt]
\item \textbf{``데이터는 원인을 말하지 않는다. 우리가 원인을 '추정'할 뿐이다.''}
\item \textbf{``AI는 집중도를 판단하지 않는다. 우리가 만든 기준을 계산할 뿐이다.''}
\item \textbf{``이 수업은 코딩을 가르치지 않는다. 코딩을 의심하는 방법을 가르친다.''}
\end{enumerate}

\vspace{3pt}
\noindent\rule{\textwidth}{0.3pt}

% ========================================
\section{2022 개정 교육과정 연계}
% ========================================

\subsection{관련 성취기준}

\begin{table}[H]
\centering
\begin{tabular}{|C{2cm}|L{9cm}|C{2.5cm}|}
\hline
\rowcolor{headerblue}
\textbf{성취기준} & \textbf{내용} & \textbf{관련 주차} \\
\hline
\textbf{[6실05-04]} & 디지털 데이터와 아날로그 데이터의 특징을 이해하고, 인공지능에 활용할 수 있는 데이터의 유형이나 형태를 탐색한다. & 1-4, 9-13주차 \\
\hline
\textbf{[6실05-05]} & 인공지능이 만들어지는 과정을 체험하고, 인공지능이 사회에 미치는 영향을 탐색한다. & 9-13주차 \\
\hline
\end{tabular}
\end{table}

\subsection{예비교사 필수 내용 지식(CK)}

\begin{table}[H]
\centering
\begin{tabular}{|L{3cm}|L{10.5cm}|}
\hline
\rowcolor{headerblue}
\textbf{CK 영역} & \textbf{세부 내용} \\
\hline
\textbf{피지컬 컴퓨팅} & 마이크로비트 구조, Grove 센서 연결, MakeCode 프로그래밍 \\
\hline
\textbf{데이터 수집} & 센서 데이터 유형, 수집 방법, 데이터 정리 \\
\hline
\textbf{데이터 분석} & 시각화, 상관관계 vs 인과관계, 반례 탐색 \\
\hline
\textbf{AI 모델} & 규칙 기반 AI, 모델의 한계와 비판적 분석 \\
\hline
\textbf{바이브 코딩} & AI 협업 코딩, 프롬프트 작성, 결과 검토 \\
\hline
\end{tabular}
\end{table}

\vspace{3pt}
\noindent\rule{\textwidth}{0.3pt}

% ========================================
\section{학습 목표}
% ========================================

\begin{itemize}[leftmargin=*, itemsep=3pt]
\item 마이크로비트와 Grove 확장 센서를 활용하여 데이터를 수집할 수 있다.
\item 수집한 데이터를 정리하고 시각화하여 패턴을 탐색할 수 있다.
\item 상관관계와 인과관계의 차이를 설명하고, 반례를 통해 가설을 검증할 수 있다.
\item 바이브 코딩을 활용하여 간단한 규칙 기반 AI 모델을 만들 수 있다.
\item AI 모델의 작동 원리와 한계를 비판적으로 분석할 수 있다.
\item 학교 현장에 적용 가능한 데이터 수집·분석 프로젝트를 기획하고 수행할 수 있다.
\end{itemize}

\vspace{3pt}
\noindent\rule{\textwidth}{0.3pt}

% ========================================
\section{수업 운영 방식}
% ========================================

\begin{table}[H]
\centering
\begin{tabular}{|C{2cm}|C{2cm}|L{9.5cm}|}
\hline
\rowcolor{headerblue}
\textbf{구분} & \textbf{주차} & \textbf{운영 방식} \\
\hline
대면 수업 & 1-4주 & 강의실 대면 수업 \\
\hline
\rowcolor{lightorange}
\textbf{데이터 수집} & \textbf{5-8주} & \textbf{실습학교에서 데이터 수집 (교육실습 기간)} \\
\hline
대면 수업 & 9-15주 & 강의실 대면 수업 \\
\hline
\end{tabular}
\end{table}

\note{실습기간(5-8주)은 교육실습 기간으로, 학생들이 실습학교에서 프로젝트 데이터를 수집합니다.}

\vspace{3pt}
\noindent\rule{\textwidth}{0.3pt}

% ========================================
\section{주차별 수업 개요}
% ========================================

\begin{longtable}{|C{1cm}|L{3cm}|L{5cm}|C{1.5cm}|L{3cm}|}
\hline
\rowcolor{headerblue}
\textbf{주차} & \textbf{주제} & \textbf{주요 내용} & \textbf{VC} & \textbf{평가} \\
\hline
\endfirsthead
\hline
\rowcolor{headerblue}
\textbf{주차} & \textbf{주제} & \textbf{주요 내용} & \textbf{VC} & \textbf{평가} \\
\hline
\endhead
\textbf{1} & 오리엔테이션 & 바이브코딩 시연, 마이크로비트 체험 & — & — \\
\hline
\textbf{2} & 바이브코딩 소개 & Claude 사용법, 센서 탐색 & O & 1주차 과제 점검 \\
\hline
\textbf{3} & 프로젝트 기획 & 주제 선정, 데이터 수집 계획서 & — & — \\
\hline
\textbf{4} & 데이터수집기 제작 & MakeCode 프로그래밍, Grove 연결 & — & — \\
\hline
\rowcolor{lightorange}
\textbf{5} & 데이터 수집 & 실습학교에서 데이터 수집 & — & — \\
\hline
\rowcolor{lightorange}
\textbf{6} & 데이터 수집 & 실습학교에서 데이터 수집 & — & — \\
\hline
\rowcolor{lightorange}
\textbf{7} & 데이터 수집 & 실습학교에서 데이터 수집 & — & — \\
\hline
\rowcolor{lightorange}
\textbf{8} & 데이터 정리 & 수집 완료, 엑셀 정리 & — & — \\
\hline
\textbf{9} & 바이브코딩 시각화 & 데이터 그래프 시각화 & O & — \\
\hline
\textbf{10} & 바이브코딩 분석 & 상관관계 vs 인과관계 & O & — \\
\hline
\textbf{11} & 바이브코딩 분석 심화 & 반례 탐색, 분석 결과 정리 & O & — \\
\hline
\textbf{12} & 바이브코딩 AI 모델 & 규칙 기반 AI 모델 구현 & O & — \\
\hline
\textbf{13} & 바이브코딩 분석 완료 & 분석 마무리, 보고서 초안 & O & — \\
\hline
\textbf{14} & 보고서 + 발표 준비 & 보고서 완성, 슬라이드 제작 & — & — \\
\hline
\textbf{15} & 최종 발표 & 발표, 피어평가, 성찰 & — & 피어평가 \\
\hline
\caption{주차별 수업 개요}
\end{longtable}

\vspace{3pt}
\noindent\rule{\textwidth}{0.3pt}

% ========================================
\section{평가 방법}
% ========================================

\subsection{평가 체계 총괄}

\begin{table}[H]
\centering
\begin{tabular}{|L{4cm}|C{2cm}|C{4cm}|L{4cm}|}
\hline
\rowcolor{headerblue}
\textbf{평가 영역} & \textbf{배점} & \textbf{평가 주차} & \textbf{세부 내용} \\
\hline
\textbf{주간 과제} & \textbf{35\%} & 2-4주, 8-13주 & 주당 5\% × 7주 \\
\hline
\textbf{피어 평가} & \textbf{5\%} & 15주차 & 동료 발표 평가 \\
\hline
\textbf{최종 프로젝트} & \textbf{60\%} & 14-15주차 & 보고서 30\% + 발표 10\% + 성찰 15\% + 피어 5\% \\
\hline
\textbf{합계} & \textbf{100\%} & & \\
\hline
\end{tabular}
\end{table}

\note{\textbf{상대평가}: A (25\%), B (50\%), C (25\%)}

\note{\textbf{1주차 과제}는 평가 점검용으로 학기 평가에 반영되지 않음}

\subsection{주차별 과제 및 제출물}

\begin{longtable}{|C{1cm}|L{2.5cm}|C{1cm}|L{5.5cm}|C{1.5cm}|}
\hline
\rowcolor{headerblue}
\textbf{주차} & \textbf{과제} & \textbf{VC} & \textbf{제출물} & \textbf{배점} \\
\hline
\endfirsthead
\hline
\rowcolor{headerblue}
\textbf{주차} & \textbf{과제} & \textbf{VC} & \textbf{제출물} & \textbf{배점} \\
\hline
\endhead
1 & LED 이름 표시 & & 사진 + 개발일지(PDF) & \textcolor{warningred}{점검용} \\
\hline
2 & 바이브코딩 체험 & O & 결과물(HTML) + 대화목록 전문(PDF) + 요약본(PDF) & 5\% \\
\hline
3 & 프로젝트 기획 & & 데이터 수집 계획서 + 개발일지(PDF) & 5\% \\
\hline
4 & 데이터수집기 제작 & & MakeCode 파일 + 테스트 결과 + 개발일지(PDF) & 5\% \\
\hline
\rowcolor{lightorange}
5-7 & 데이터 수집 & & 없음 (8주차에 통합 제출) & - \\
\hline
\rowcolor{lightorange}
8 & 데이터 정리 & & 엑셀 파일 + 개발일지(PDF) & 5\% \\
\hline
9 & 바이브코딩 시각화 & O & 결과물(HTML) + 대화목록 전문(PDF) + 요약본(PDF) & 5\% \\
\hline
10 & 바이브코딩 분석 & O & 결과물(HTML) + 해석 메모 + 대화목록(PDF) & 5\% \\
\hline
11 & 바이브코딩 분석 심화 & O & 결과물(HTML) + 반례 분석 메모 + 대화목록(PDF) & 5\% \\
\hline
12 & 바이브코딩 AI 모델 & O & 결과물(HTML) + 대화목록(PDF) & 5\% \\
\hline
13 & 바이브코딩 분석 완료 & O & 결과물(HTML) + 보고서 초안(PDF) + 대화목록(PDF) & 5\% \\
\hline
\caption{주차별 과제 및 제출물 (VC = 바이브코딩)}
\end{longtable}

\note{\textbf{VC = 바이브코딩}. O 표시 주차는 대화목록 전문 + 요약본 필수 제출}

\note{주황색 배경: 실습기간 (실습학교에서 데이터 수집)}

\subsection{주간 과제 등급 기준}

\begin{table}[H]
\centering
\begin{tabular}{|C{1.5cm}|C{1.5cm}|L{10cm}|}
\hline
\rowcolor{headerblue}
\textbf{등급} & \textbf{점수} & \textbf{기준} \\
\hline
A & 100\% & 과제 완벽 수행 + 개발일지 충실 + 추가 시도/통찰 \\
\hline
B & 85\% & 과제 대부분 수행 + 개발일지 충실 \\
\hline
C & 70\% & 과제 일부 수행 + 개발일지 기본 작성 \\
\hline
D & 55\% & 과제 미흡 또는 개발일지 부실 \\
\hline
F & 0\% & 미제출 \\
\hline
\end{tabular}
\end{table}

\subsection{피어 평가 (5\%)}

\begin{table}[H]
\centering
\begin{tabular}{|L{3cm}|L{10.5cm}|}
\hline
\rowcolor{headerblue}
\textbf{항목} & \textbf{내용} \\
\hline
\textbf{실시 시기} & 15주차 최종 발표 시 \\
\hline
\textbf{평가 대상} & 동료의 발표 \\
\hline
\textbf{평가 기준} & 발표 내용, 통찰력, 완성도 \\
\hline
\end{tabular}
\end{table}

\begin{table}[H]
\centering
\begin{tabular}{|C{1.5cm}|L{12cm}|}
\hline
\rowcolor{headerblue}
\textbf{점수} & \textbf{기준} \\
\hline
3점 & 우수 - 통찰력 있는 분석, 창의적 시도 \\
\hline
2점 & 보통 - 기본 요구사항 충족 \\
\hline
1점 & 미흡 - 내용 일부만 충족 \\
\hline
0점 & 부족 - 발표 불참 또는 미완성 \\
\hline
\end{tabular}
\end{table}

\subsection{최종 프로젝트 (60\%)}

\begin{table}[H]
\centering
\begin{tabular}{|L{4cm}|C{1.5cm}|L{8cm}|}
\hline
\rowcolor{headerblue}
\textbf{구성 요소} & \textbf{배점} & \textbf{핵심 평가 기준} \\
\hline
\textbf{프로젝트 보고서} & 30\% & 프로젝트 과정 명확 + 데이터 분석 적절 + 통찰 구체적 \\
\hline
\textbf{발표} & 10\% & 발표 내용 명확 + 전달력 우수 \\
\hline
\textbf{개인 성찰 보고서} & 15\% & 학습 목표 달성 명확 + 관점 변화 구체적 + 학교 적용 방안 \\
\hline
\textbf{피어 평가} & 5\% & 동료 발표 평가 \\
\hline
\textbf{합계} & \textbf{60\%} & \\
\hline
\end{tabular}
\end{table}

\vspace{3pt}
\noindent\rule{\textwidth}{0.3pt}

% ========================================
\section{과제 제출 방법}
% ========================================

\subsection{바이브코딩 과제 제출물}

바이브코딩 활용 주차(2, 9-13주)에는 \textbf{대화목록 전문}과 \textbf{대화목록 요약본}을 필수 제출한다.

\begin{table}[H]
\centering
\begin{tabular}{|L{4cm}|L{9.5cm}|}
\hline
\rowcolor{headerblue}
\textbf{제출물} & \textbf{내용} \\
\hline
\textbf{대화목록 전문} & 프롬프트 원문 + Claude 응답 전체 \\
\hline
\textbf{대화목록 요약본} & 프롬프트 원문 + Claude 응답 3줄 요약 \\
\hline
\end{tabular}
\end{table}

\subsection{대화목록 PDF 생성 방법}

대화목록은 Claude에게 아래 프롬프트를 입력하면 자동으로 PDF가 생성된다.

\begin{tcolorbox}[colback=lightorange, colframe=accentorange, boxrule=0.5pt, title={\textbf{대화목록 PDF 생성 프롬프트}}, fonttitle=\bfseries\color{accentorange!80!black}]
\small
이 대화에서 사용한 프롬프트와 답변을 포함한 대화 전문과 프롬프트와 답변을 정리한 요약본을 PDF 문서로 만들어줘. 오늘 날짜, 수업주차 O주차, 작성자 이름 OOO, 학번 XXXXXX을 추가해줘.
\end{tcolorbox}

\textbf{사용 예시:}
\begin{quote}
\small
이 대화에서 사용한 프롬프트와 답변을 포함한 대화 전문과 프롬프트와 답변을 정리한 요약본을 PDF 문서로 만들어줘. 오늘 날짜, 수업주차 \textbf{2주차}, 작성자 이름 \textbf{홍길동}, 학번 \textbf{202012345}을 추가해줘.
\end{quote}

\subsection{제출 안내}

\begin{table}[H]
\centering
\begin{tabular}{|L{3cm}|L{10.5cm}|}
\hline
\rowcolor{headerblue}
\textbf{항목} & \textbf{내용} \\
\hline
\textbf{제출처} & LMS \\
\hline
\textbf{제출 기한} & 매주 일요일 자정 \\
\hline
\textbf{지각 감점} & 1일당 10\% 감점 \\
\hline
\textbf{출석 인정} & 과제 제출로 출석 대체 \\
\hline
\end{tabular}
\end{table}

\vspace{3pt}
\noindent\rule{\textwidth}{0.3pt}

% ========================================
\section{수업 자료 및 도구}
% ========================================

\subsection{필수 준비물}

\begin{table}[H]
\centering
\begin{tabular}{|L{5cm}|L{8.5cm}|}
\hline
\rowcolor{headerblue}
\textbf{항목} & \textbf{비고} \\
\hline
마이크로비트 V2 & 학교 대여 \\
\hline
Grove 확장 보드 + 센서 & 학교 대여 \\
\hline
USB 케이블 & 학교 대여 \\
\hline
노트북/PC & 개인 지참 \\
\hline
Claude 계정 & 학교 제공 예정 \\
\hline
\end{tabular}
\end{table}

\subsection{소프트웨어}

\begin{table}[H]
\centering
\begin{tabular}{|L{3.5cm}|L{5cm}|L{5cm}|}
\hline
\rowcolor{headerblue}
\textbf{도구} & \textbf{용도} & \textbf{URL} \\
\hline
MakeCode & 마이크로비트 프로그래밍 & makecode.microbit.org \\
\hline
Claude & 바이브코딩 & claude.ai \\
\hline
엑셀/구글 시트 & 데이터 정리 & — \\
\hline
\end{tabular}
\end{table}

\subsection{학습 윤리}

\begin{table}[H]
\centering
\begin{tabular}{|L{6.5cm}|L{6.5cm}|}
\hline
\rowcolor{lightgreen}
\textbf{권장} & \textbf{금지} \\
\hline
AI 코드를 수정·검토 & AI 코드 무검증 복붙 \\
\hline
``이런 느낌으로'' 자연어 설명 & ``AI가 이렇게 했으니까 맞다'' \\
\hline
AI 코드의 문제점 지적 & AI를 정답 생성기로 사용 \\
\hline
\end{tabular}
\end{table}

\vspace{3pt}
\noindent\rule{\textwidth}{0.3pt}

% ========================================
\section*{참고자료}
\addcontentsline{toc}{section}{참고자료}
% ========================================

\begin{table}[H]
\centering
\begin{tabular}{|C{1cm}|L{6cm}|L{2.5cm}|L{4cm}|}
\hline
\rowcolor{headerblue}
\textbf{\#} & \textbf{자료명} & \textbf{유형} & \textbf{비고} \\
\hline
1 & 문서버전관리가이드\_v1.2.2 & 문서 & 버전 관리 규칙 \\
\hline
2 & 마이크로비트 공식 문서 & 웹사이트 & microbit.org \\
\hline
3 & Grove 센서 문서 & 웹사이트 & wiki.seeedstudio.com \\
\hline
4 & 2022 개정 교육과정 & 교육부 문서 & 실과, SW/AI 교육 \\
\hline
\end{tabular}
\end{table}

\vspace{3pt}
\noindent\rule{\textwidth}{0.3pt}

% ========================================
\section*{생성/수정 이력}
\addcontentsline{toc}{section}{생성/수정 이력}
% ========================================

\begin{table}[H]
\centering
\small
\begin{tabular}{|L{1.8cm}|L{2cm}|L{1.5cm}|L{8cm}|}
\hline
\rowcolor{headerblue}
\textbf{버전} & \textbf{날짜} & \textbf{수준} & \textbf{변경 내용} \\
\hline
v1.0.0.0 & 2026-01-24 & 최초 & 초안 작성 \\
\hline
v1.5.0.0 & 2026-01-26 & 마이너 & 프로젝트 중심 재구성: 학생 주도 프로젝트, 실습기간 데이터 수집, 바이브코딩 소개 추가 \\
\hline
v1.5.1.0 & 2026-01-26 & 리비전 & 9-13주 바이브코딩 활용 명시, 대화목록 필수 제출로 통일 \\
\hline
v1.5.2.0 & 2026-01-26 & 리비전 & 1주차 과제 평가 점검용 명시, 배점 조정(주간과제 35\%, 최종프로젝트 60\%) \\
\hline
v1.5.3.0 & 2026-01-26 & 리비전 & 제출물 명칭 변경: 대화목록 전문 + 요약본 \\
\hline
\textbf{v1.5.4.0} & \textbf{2026-02-13} & \textbf{리비전} & \textbf{LaTeX 교재형식 변환, AI디지털리터러시 스타일 통일, 주차별 수업 개요 위치 이동(7장)} \\
\hline
\end{tabular}
\end{table}

\vspace{15pt}
\begin{center}
\textit{사물인터넷과 빅데이터 — 15주 강의계획서}\\[3pt]
\textit{청주교육대학교 컴퓨터교육과}
\end{center}

\end{document}
